\begin{pastexam}[label={px:nccu-110-dsa-03}]{政大資科所 110 學年度 第 3 題}
（問答）

(15\%) Use an example to explain why Dijkstra's shortest path algorithm cannot be applied to graphs
with negative edge weights. In your example graph, please highlight the source and the destination.
\end{pastexam}

\begin{pastsoln}
\textbf{例子}。 有向圖，起點 $s$（來源），終點 $t$（目的地），四條邊：$s \to a$ 權重 $2$、$s \to b$ 權重 $3$、
$b \to a$ 權重 $-2$、$a \to t$ 權重 $2$。 圖中雙圈的是起點和終點。
\begin{center}
\begin{tikzpicture}[x=1.6cm,y=1.2cm,>={Stealth[length=2.4mm]},
  v/.style={circle,draw,minimum size=7mm,inner sep=0pt},
  st/.style={circle,draw,double,very thick,minimum size=7mm,inner sep=0pt},
  w/.style={fill=white,inner sep=1.5pt,font=\small}]
\node[st] (s) at (0,0) {$s$}; \node[font=\small,left=2pt of s] {起點};
\node[v] (a) at (1.5,0.9) {$a$};
\node[v] (b) at (1.5,-0.9) {$b$};
\node[st] (t) at (3,0.9) {$t$}; \node[font=\small,right=2pt of t] {終點};
\draw[->] (s) -- node[w] {$2$} (a);
\draw[->] (s) -- node[w] {$3$} (b);
\draw[->] (b) -- node[w] {$-2$} (a);
\draw[->] (a) -- node[w] {$2$} (t);
\end{tikzpicture}
\end{center}
真正的最短路徑是 $s \to b \to a \to t$，長度 $3 + (-2) + 2 = 3$；另一條 $s \to a \to t$ 長度 $4$。 圖裡沒有負環，
最短距離是有定義的：$\delta(s, a) = 1$、$\delta(s, b) = 3$、$\delta(s, t) = 3$。

\textbf{Dijkstra 的執行}（資結書虛擬碼~25.3）。 每一步取出 $d$ 最小、還沒確定的頂點，放進 $S$（認定它的 $d$ 已經是
最短距離），再鬆弛它的出邊：
\begin{center}
\begin{tabular}{clcccc}
\toprule
步驟 & 取出 & $d[s]$ & $d[a]$ & $d[b]$ & $d[t]$ \\
\midrule
初始 & & $0$ & $\infty$ & $\infty$ & $\infty$ \\
1 & $s$；鬆弛 $s \to a$、$s \to b$ & $0$ & $2$ & $3$ & $\infty$ \\
2 & $a$（$d = 2$ 最小）；鬆弛 $a \to t$ & $0$ & $\mathbf{2}$ & $3$ & $4$ \\
3 & $b$；$a$ 已經在 $S$，不再更新 & $0$ & $2$ & $3$ & $4$ \\
4 & $t$ & $0$ & $2$ & $3$ & $\mathbf{4}$ \\
\bottomrule
\end{tabular}
\end{center}
Dijkstra 回答 $s$ 到 $t$ 的距離是 $4$、路徑 $s \to a \to t$，但正確答案是 $3$。 $d[a] = 2$ 也錯了（應該是 $1$）。

\textbf{錯在哪裡}。 Dijkstra 的正確性證明（資結書定理~25.6）靠的是：取出 $u$ 時，任何還沒確定的頂點 $y$ 都有
$d[y] \ge d[u]$，而從 $y$ 再走下去，權重非負，路徑只會變長，所以不可能經過 $y$ 繞回來得到更短的 $u$。 有負邊時
這個推論不成立：第 2 步取出 $a$ 時 $d[b] = 3 > 2$，看起來經過 $b$ 只會更遠，但 $b \to a$ 的 $-2$ 讓 $3 - 2 = 1 < 2$。
$a$ 一旦放進 $S$ 就不會再被修正，靠 $a$ 算出的 $d[t]$ 也跟著錯。

就算實作時允許第 3 步把 $d[a]$ 改成 $1$（不檢查 $a$ 是否在 $S$ 裡），$a$ 已經取出過，不會再去鬆弛 $a \to t$，
$d[t]$ 還是 $4$，一樣錯。 有負邊（沒有負環）時要改用 Bellman--Ford（資結書第~25.3~節），它對這張圖算出
$d[t] = 3$；圖是 DAG 時也可以用拓樸順序鬆弛（資結書第~25.5~節）。

補充：題目若要無向圖的例子，要小心一條負邊 $\{u, v\}$ 在無向圖裡就等於 $u \to v \to u$ 這個負環，最短路徑沒有定義，
所以這類例子通常用有向圖舉。
\end{pastsoln}
